\documentclass[11pt]{article}

\usepackage[margin=1in]{geometry}
\usepackage{amsmath, amssymb, amsthm, bm, mathtools}
\usepackage{booktabs}
\usepackage{array}
\usepackage{enumitem}
\usepackage{hyperref}
\usepackage{xcolor}
\usepackage{microtype}

\hypersetup{
  hypertexnames=false,
  colorlinks=true,
  linkcolor=blue!45!black,
  citecolor=blue!45!black,
  urlcolor=blue!45!black
}

\urlstyle{same}
\allowdisplaybreaks
\setlist[itemize]{leftmargin=1.8em}
\setlist[enumerate]{leftmargin=1.8em}
\emergencystretch=3em

\newcommand{\E}{\mathbb{E}}
\newcommand{\Var}{\operatorname{Var}}
\newcommand{\Cov}{\operatorname{Cov}}
\newcommand{\Skew}{\operatorname{Skew}}
\newcommand{\Kurt}{\operatorname{Kurt}}
\newcommand{\kcum}{\kappa}
\newcommand{\tr}{\operatorname{tr}}
\newcommand{\vecop}{\operatorname{vec}}
\newcommand{\R}{\mathbb{R}}
\newcommand{\dd}{\mathrm{d}}
\newcommand{\mc}[1]{\mathcal{#1}}
\newcommand{\mb}[1]{\mathbf{#1}}
\newcommand{\bs}[1]{\boldsymbol{#1}}
\newcommand{\Oh}{\Omega_t}
\newcommand{\Obar}{\bar\Omega}
\newcommand{\DelOh}{\Delta\Omega_t}
\newcommand{\taux}{\tau_{x,t}}
\newcommand{\taul}{\tau_{\ell,t}}
\newcommand{\Lamell}{\Lambda_{\ell,t}}
\newcommand{\sigA}{\sigma_{A,t}}
\newcommand{\sige}{\sigma_{e,t}}
\newcommand{\sigd}{\sigma_{d,t}}
\newcommand{\rk}{r^k_{t+1}}
\newcommand{\cE}{c^E_t}
\newcommand{\cW}{c^W_t}
\newcommand{\cEp}{c^E_{t+1}}
\newcommand{\cWp}{c^W_{t+1}}
\newcommand{\sE}{s^E_{t+1}}
\newcommand{\sW}{s^W_{t+1}}
\newcommand{\Dthree}{\mathcal{D}^3}
\newcommand{\zhat}{\hat z_t}
\newcommand{\Dz}{\Delta z_t}
\newcommand{\kxtwo}{\kappa_x^{(2)}}
\newcommand{\kxthree}{\kappa_x^{(3)}}
\newcommand{\kxvov}{\kappa_x^{(3,\mathrm{vov})}}
\newcommand{\kxcross}{\kappa_x^{(3,\mathrm{cross})}}
\newcommand{\nuth}{\nu_\theta}
\newcommand{\nuthIII}{\nu_\theta^{(3)}}
\newcommand{\nuthx}{\nu_\theta^{(\mathrm{cross})}}

\theoremstyle{plain}
\newtheorem{theorem}{Theorem}
\newtheorem{proposition}{Proposition}
\newtheorem{lemma}{Lemma}

\theoremstyle{definition}
\newtheorem{assumption}{Assumption}
\newtheorem{definition}{Definition}
\newtheorem{remark}{Remark}

\title{Second-Moment Wedge Projection and the FVGQ Risk-Sharing Channel\\
\large A Standalone Research Report}
\author{Caught in the Crossfire Research Notes}
\date{May 1, 2026}

\begin{document}

\maketitle

\begin{abstract}
This report turns the informal argument that ``a first-order BCA prototype
projects second-moment uncertainty into first-order wedges'' into a
self-contained theoretical document. We work entirely within the
Fern\'andez-Villaverde--Guerr\'on-Quintana (2020; FVGQ) two-agent collateral
RBC, whose three-stage timing, skin-in-the-game financial friction, and three
stochastic-volatility shocks are reproduced in §3 so that the report can be
read without recourse to the original paper. The central second-order result
is a wedge projection equation: when the true data-generating process contains
stochastic volatility, the investment wedge recovered by a Business Cycle
Accounting (BCA) prototype must satisfy a forward-looking equation forced by
the second-order risk correction in the true Euler equation; under AR(1)
volatility, the projection admits a closed form, decomposable into
precautionary, risk-premium, collateral-risk, and two-agent aggregation
components. We then advance the projection to third order in the perturbation
parameter (Theorem~\ref{thm:third-order}). The third-order extension
introduces a vol-of-vol refinement of the linear coefficient, a state-dependent
cross term $\Gamma(\hat z_t,\Omega_t)$, and an explicit asymmetric response of
the BCA wedge to positive versus negative volatility innovations. These three
new objects supply identifying handles that the second-order theory cannot
deliver, including a structural explanation for the negative wide-grid Monte
Carlo estimate of the collateral-risk sensitivity $\nu_\theta$. The report
also contrasts the FVGQ collateral channel with a first-order working-capital
channel, which projects into the after-tax labor wedge rather than the
investment wedge.
\end{abstract}

\tableofcontents
\clearpage

\section{Purpose and Main Results}

This report integrates the following six internal research notes:
\begin{itemize}
  \item \texttt{2026-05-01-wedge-projection-derivation-skeleton.md}
  \item \texttt{2026-05-01-fvgq-Phi-closed-form.md}
  \item \texttt{2026-05-01-fvgq-Phi-numerical-validation.md}
  \item \texttt{2026-05-01-fvgq-two-agent-risk-sharing-wedge.md}
  \item \texttt{2026-05-01-fvgq-nu-theta-monte-carlo.md}
  \item \texttt{2026-05-01-fvgq-aggregation-wedge-numerical.md}
\end{itemize}
The common question is how a first-order BCA prototype absorbs the
second-order uncertainty terms generated by a nonlinear DSGE model with
stochastic volatility.

The main conclusions are as follows.
\begin{enumerate}
  \item The BCA investment wedge satisfies the projection equation
  \[
  \hat\tau_{x,t}
  =
  \omega_\tau \E_t\hat\tau_{x,t+1}
  +\frac{1}{2}\Phi(\Omega_t),
  \]
  where $\Phi(\Omega_t)$ is the second-order risk correction in the true Euler
  equation.

  \item If $\Omega_t$ follows an AR(1) process, the projection has the closed
  form
  \[
  \hat\tau_{x,t}
  =
  \frac{\Phi(\Obar)}{2(1-\omega_\tau)}
  +
  \frac{\phi}{2(1-\omega_\tau\rho_\Omega)}
  (\Omega_t-\Obar),
  \qquad
  \phi=\Phi'(\Obar).
  \]
  Thus BCA can transform a second-moment stochastic-volatility signal into a
  persistent first-order investment wedge.

  \item In an FVGQ-style financial-friction model, the risk correction can be
  decomposed as
  \[
  \Phi^{\mathrm{FVGQ}}
  =
  \Phi^{\mathrm{prec}}
  +
  \Phi^{\mathrm{rp}}
  +
  \Phi^{\mathrm{coll}}
  +
  \Phi^{\mathrm{agg}}.
  \]
  These terms capture precautionary motives, risk premia, collateral risk, and
  the aggregation error induced by using total consumption in a two-agent
  economy.

  \item Financial uncertainty in the FVGQ collateral mechanism is predicted to
  appear mainly in the investment wedge, with a positive sign under the
  baseline collateral-risk channel. The sign of macro uncertainty is not
  guaranteed, because the risk-premium covariance can dominate the
  precautionary component.

  \item Numerical diagnostics from the FVGQ workspace qualify the simple
  closed-form comparison. The Monte Carlo estimate of the collateral-risk term
  $\nu_\theta$ is small and negative in the current grid, while the two-agent
  aggregation correction is positive for both volatility channels. Aggregation
  therefore cannot explain the negative local-projection peak for the macro
  uncertainty shock.

  \item A working-capital mechanism has a different empirical signature. If
  firms must finance wage payments in advance, the BCA after-tax labor wedge
  satisfies
  \[
  \hat\Lambda_{\ell,t}
  =
  -\lambda_{wc}(q_t-\bar q)+O((q_t-\bar q)^2),
  \qquad
  \lambda_{wc}>0.
  \]
  Therefore a rise in the external finance spread lowers the after-tax labor
  wedge. This sign is first-order and does not depend on Hessian terms.

  \item Extending the projection equation to third order in the perturbation
  parameter (Theorem~\ref{thm:third-order}) introduces three new objects: a
  refinement to the second-order coefficient via vol-of-vol amplification, a
  state-dependent cross term $\Gamma(\hat z_t,\Omega_t)$, and an asymmetric
  response of $\hat\tau_{x,t}$ to positive versus negative volatility
  innovations. The third-order extension is necessary to interpret the FVGQ
  generalized impulse response functions, because in second order alone the
  on-impact response to an uncertainty shock $u_{j,t}$ is identically zero;
  it is only at third order that $u_{j,t}$ enters the policy function.
\end{enumerate}

\section{Notation and Normalization}

Hats denote log deviations from the deterministic steady state; for example,
$\hat C_t=\log(C_t/\bar C)$. The stochastic-volatility state is
\[
\Omega_t=(\sigma_{A,t}^2,\sigma_{e,t}^2)' ,
\qquad
\Delta\Omega_t=\Omega_t-\bar\Omega .
\]
In the one-factor case, $\Omega_t$ can be treated as a scalar.

To avoid ambiguity in the labor-wedge sign convention, this report uses the
after-tax labor wedge
\[
\Lambda_{\ell,t}\equiv 1-\tau_{\ell,t}^{\mathrm{tax}}.
\]
Thus $\hat\Lambda_{\ell,t}<0$ means that the labor wedge deteriorates. If the
same result is reported in terms of the labor-tax rate
$\tau_{\ell,t}^{\mathrm{tax}}$, the sign is reversed.

\begin{remark}[Normalization of $\Phi$]
This report uses the convention
\[
\hat\tau_{x,t}
=
\omega_\tau\E_t\hat\tau_{x,t+1}
+\frac{1}{2}\Phi(\Omega_t).
\]
Some of the underlying notes absorb the factor $1/2$ into the definition of
$\Phi$. All formulas for $\phi_A$, $\phi_e$, $\kappa_{x,A}$, and
$\kappa_{x,e}$ in this report use the convention above. A different convention
rescales the coefficients but leaves the sign-identification logic unchanged.
\end{remark}

\section{The FVGQ 2020 Model Environment}
\label{sec:fvgq-model}

The closed-form risk decomposition in §\ref{sec:fvgq-closed-form}, the two-agent
aggregation correction in §\ref{sec:two-agent}, and the numerical validation in
§\ref{sec:numerical-validation} all take Fern\'andez-Villaverde and
Guerr\'on-Quintana (2020; henceforth FVGQ) as the reference
data-generating process. The original paper presents the model in its Section~5.
This section reproduces the equations, conventions, and calibration needed for
the projection theory below, so that the report is self-contained.

\subsection{Households, Three-Stage Timing, and Equity Claims}

Time is discrete. There is a continuum of ex-ante identical households of
measure one, each containing a continuum of members. Aggregate state is
Markov, and a representative household head writes contingency plans for all
members.

\medskip

\noindent\textbf{Three stages.} Each period is divided into three stages.
\begin{enumerate}
  \item \textit{Decision stage.} Members are united and share assets, which
  consist of $s_t$ units of an equity claim with price $q_t$. One claim
  entitles the holder to one unit of capital next period. The household head
  splits assets across members and writes contingent allocation plans.

  \item \textit{Production stage.} Each member draws a shock determining her
  type. With probability $\pi$ the member becomes an entrepreneur; with
  probability $1-\pi$ she becomes a worker. Entrepreneurs invest in capital
  but do not work; workers supply labor but cannot invest. Both can trade
  equity claims, subject to the financial frictions specified below.

  \item \textit{Consumption-investment stage.} Workers receive wage income;
  entrepreneurs collect capital rental income, depreciate their equity, and
  invest $i_t$ units of the consumption good through a one-to-one technology
  yielding $i_t$ units of new capital. Consumption takes place. The period
  ends; member identities are reshuffled.
\end{enumerate}

The fraction $\pi$ is calibrated to $\pi=0.06$, matching the share of firms
that invest in any given quarter in U.S.\ Compustat data.

\subsection{Production Technology and Capital Utilization}

A representative competitive firm produces a final good using capital services
$k_t^D$ and labor $\ell_t$:
\begin{equation}
y_t = A_t (k_t^D)^\alpha \ell_t^{1-\alpha},
\qquad
k_t^D = u_t k_t,
\tag{Y-FVGQ}
\end{equation}
where $u_t$ is the capital utilization rate. The utilization rate controls
depreciation through the reduced-form schedule
\begin{equation}
\delta(u_t) = \delta_0 + \delta_1 (u_t-1) + \frac{\delta_2}{2} (u_t-1)^2,
\tag{Dep-FVGQ}
\end{equation}
with $\delta_0$ the deterministic SS depreciation rate. A higher $u_t$
accelerates depreciation, providing the firm with an internal margin to absorb
shocks without changing $k_t$.

TFP follows a level AR(1) with stochastic volatility:
\begin{equation}
A_t = (1-\rho_A) \bar A + \rho_A A_{t-1} + e^{\sigA} \varepsilon_{A,t},
\qquad
\varepsilon_{A,t} \sim \mc{N}(0,1),
\qquad
\bar A = 1.
\tag{A-FVGQ}
\end{equation}
The reflecting barrier $\bar A>0$ guarantees positivity but is irrelevant for
local approximations. The conditional volatility $\sigA$ is itself a stochastic
process described in §\ref{subsec:fvgq-uncertainty} below.

The firm's first-order conditions equate marginal products to factor prices:
\[
\alpha \frac{y_t}{k_t^D} = r_t,
\qquad
(1-\alpha) \frac{y_t}{\ell_t} = w_t.
\]

\subsection{Skin-in-the-Game Financial Friction}

The defining feature of FVGQ is a Kiyotaki-Moore-style collateral constraint
on the entrepreneur's portfolio. Two parameters control how much of the
entrepreneur's investment can be financed externally.
\begin{itemize}
  \item For \emph{new capital}: the entrepreneur can sell at most a fraction
  $\theta_t \in [0,1]$ of new equity to outside investors, retaining at least
  $(1-\theta_t) i_t$. The pledgeability $\theta_t$ is time-varying.

  \item For \emph{existing capital} $(1-\delta(u_t)) s_t$: the entrepreneur
  can sell at most a fraction $\phi \in [0,1]$ to outside investors, retaining
  $(1-\phi) (1-\delta(u_t)) s_t$. The resaleability $\phi$ is constant.
\end{itemize}

Combining the two restrictions delivers the skin-in-the-game lower bound on
entrepreneur equity holdings:
\begin{equation}
\sE \;\ge\; (1-\theta_t) i_t + (1-\phi) (1-\delta(u_t)) s_t.
\tag{F-FVGQ}
\end{equation}
The analogous bound for workers,
$\sW \ge (1-\phi)(1-\delta(u_t)) s_t$, never binds in equilibrium because
workers are net buyers of equity.

This formulation differs from environments with idiosyncratic productivity
risk (Bernanke-Gertler-Gilchrist, Buera-Quadrini-Shin) in that the friction
scales with the size of investment rather than with productivity heterogeneity.
It also differs from Jermann and Quadrini (2012) in that $\theta_t$ is exogenous
and time-varying; this exogeneity is what supplies the second source of
stochastic volatility.

The constraint (F-FVGQ) binds in equilibrium whenever the price of capital
exceeds one ($q_t>1$), which is the relevant case throughout the
calibration.

\subsection{Three Uncertainty Shocks (Stochastic Volatility)}
\label{subsec:fvgq-uncertainty}

The model has three structural level shocks --- TFP $A_t$, preference shifter
$d_t$, and pledgeability $\theta_t$ --- each driven by an innovation with
time-varying conditional volatility:
\begin{align}
A_t &= (1-\rho_A) \bar A + \rho_A A_{t-1} + e^{\sigA} \varepsilon_{A,t},
\tag{Lvl-A}\\
d_t &= (1-\rho_d) \bar d + \rho_d d_{t-1} + e^{\sigd} \varepsilon_{d,t},
\tag{Lvl-d}\\
\theta_t &= (1-\rho_\theta) \bar\theta + \rho_\theta \theta_{t-1}
+ e^{\sigma_{\theta,t}} \varepsilon_{\theta,t}.
\tag{Lvl-$\theta$}
\end{align}
Each conditional volatility itself follows an AR(1):
\begin{equation}
\sigma_{j,t} = (1-\rho_{\sigma_j}) \sigma_j + \rho_{\sigma_j} \sigma_{j,t-1}
+ (1-\rho_{\sigma_j}^2)^{1/2} \nu_j u_{j,t},
\qquad
u_{j,t} \sim \mc{N}(0,1),
\tag{SV-$j$}
\end{equation}
for $j \in \{A, d, \theta\}$. The innovations $u_{j,t}$ are the
\emph{uncertainty shocks}: $u_{A,t}$ is the supply-side TFP uncertainty
shock; $u_{d,t}$ is the demand-side preference uncertainty shock; and
$u_{\theta,t}$ is the financial-friction uncertainty shock.

\medskip

\noindent\textbf{Notation reconciliation with the implementation.}
The remainder of this report (and the Dynare implementation in
\texttt{fvgq2020\_solveonly.mod}) writes the financial-friction
volatility as $\sige$ rather than $\sigma_{\theta,t}$, with associated
innovation $u_{e,t}$ rather than $u_{\theta,t}$. The two notations refer to
the same object. Sections \ref{sec:fvgq-closed-form}--\ref{sec:numerical-validation}
use $\sige$ and $u_{e,t}$ to match the codebase.

\subsection{Optimality Conditions and the Liquidity Wedge}

The household head maximizes
\begin{equation}
\E_t \sum_{t=0}^\infty \beta^t d_t
\left\{
\pi \frac{(\cE)^{1-\rho}-1}{1-\rho}
+ (1-\pi) \frac{[\cW (1-\ell_t)^\eta]^{1-\rho}-1}{1-\rho}
\right\}
\tag{U-FVGQ}
\end{equation}
subject to a household-level budget constraint and the financial constraint
(F-FVGQ). Here $\rho$ is the inverse intertemporal elasticity of substitution,
$\eta>0$ controls leisure curvature, and $d_t$ is the demand shifter described
in (Lvl-d).

After substituting the binding (F-FVGQ) into the household budget constraint,
the modified budget constraint contains a positive premium on investment
whenever $q_t>1$. The Lagrange multiplier on this modified constraint defines
a \emph{liquidity wedge}:
\begin{equation}
\lambda_t = \frac{q_t - 1}{1 - \theta_t q_t}.
\tag{Liq-FVGQ}
\end{equation}
Provided $1 < q_t < 1/\theta_t$, this wedge is positive: it prices the value
of one unit of liquidity to the household head, who can convert it into
$1/(1-\theta_t q_t)$ units of capital via leveraged investment.

The household-head optimality conditions take the following form.

\begin{itemize}
\item \textbf{Worker labor supply.} Because workers face no borrowing
constraint, the labor supply margin is undistorted by the financial friction:
\begin{equation}
\eta \frac{\cW}{1-\ell_t} = w_t.
\tag{LS-FVGQ}
\end{equation}

\item \textbf{Intra-household consumption allocation.} Entrepreneurs face a
liquidity premium, so their marginal utility exceeds that of workers:
\begin{equation}
(\cE)^{-\rho} = (1+\lambda_t) (\cW)^{-\rho} (1-\ell_t)^{\eta(1-\rho)}.
\tag{IH-FVGQ}
\end{equation}

\item \textbf{Investment Euler (capital pricing).} The price of equity claims
satisfies
\begin{equation}
q_t = \E_t \!\left\{
\beta \frac{d_{t+1}}{d_t}
\left(\frac{\cWp}{\cW}\right)^{-\rho}
\left(\frac{1-\ell_{t+1}}{1-\ell_t}\right)^{\eta(1-\rho)}
\mc{P}_{t+1}
\right\},
\tag{IE-FVGQ}
\end{equation}
where the per-claim payoff is
\begin{equation}
\mc{P}_{t+1} = r_{t+1} u_{t+1} + q_{t+1} (1-\delta(u_{t+1}))
+ \pi \lambda_{t+1} \big[r_{t+1} u_{t+1} + q_{t+1} \phi (1-\delta(u_{t+1}))\big].
\tag{Pay-FVGQ}
\end{equation}
The first two summands are the dividend and the depreciated capital gain. The
$\pi \lambda_{t+1}$ term is the value of the liquidity services that the
equity claim provides to entrepreneurs in the next period.

\item \textbf{Capital utilization.} The marginal benefit of using capital
more intensively equals the marginal cost from faster depreciation:
\begin{equation}
r_t = q_t \delta'(u_t).
\tag{Util-FVGQ}
\end{equation}

\item \textbf{Capital accumulation.}
\begin{equation}
k_{t+1} = (1-\delta(u_t)) k_t + \pi i_t.
\tag{Cap-FVGQ}
\end{equation}

\item \textbf{Goods market clearing.}
\begin{equation}
\pi \cE + (1-\pi) \cW + \pi i_t = y_t.
\tag{GM-FVGQ}
\end{equation}
\end{itemize}

\subsection{Calibration}
\label{subsec:fvgq-calibration}

Table~\ref{tab:fvgq-calib} reproduces the calibration used in
\texttt{fvgq2020\_solveonly.mod} and throughout the report. Two parameters
($\delta_1$ and $\eta$) are pinned down by SS targets; all others are taken
directly from the calibration in FVGQ Section~7.

\begin{table}[h]
\centering
\caption{FVGQ 2020 calibration used throughout the report.}
\label{tab:fvgq-calib}
\begin{tabular}{lcl}
\toprule
Parameter & Value & Target / Source \\
\midrule
$\beta$: discount factor & $0.994$ & exogenously chosen \\
$\rho$: relative risk aversion & $2$ & exogenously chosen \\
$\pi$: fraction of entrepreneurs & $0.06$ & investing-firm fraction $\simeq 0.24$/yr \\
$\eta$: leisure curvature & $1.525$ & hours of work $=1/3$ \\
$\alpha$: capital share & $0.36$ & labor income share $=0.64$ \\
$\delta_0$: SS depreciation & $0.0163$ & annual investment/capital $=0.065$ \\
$\delta_1$: linear utilization slope & $0.0131$ & pinned by SS \\
$\delta_2$: utilization curvature & $0.33\,\delta_1$ & Com\'in and Gertler (2006) \\
$\phi$: existing-capital resaleability & $0.0954$ & capital/output $=3.30$ \\
$\bar\theta$: SS new-capital pledgeability & $0.0954$ & set equal to $\phi$ \\
$\rho_A=\rho_d=\rho_\theta$: level persistences & $0.95$ & Cooley and Prescott (1995) \\
$\rho_{\sigma_A}=\rho_{\sigma_d}=\rho_{\sigma_\theta}$: SV persistences & $0.75$ & FVGQ et al.\ (2015a) \\
$\sigma_A$: SS TFP volatility & $\log(0.007)$ & Cooley and Prescott (1995) \\
$\sigma_d$: SS preference volatility & $\log(0.032)$ & FVGQ et al.\ (2015a) \\
$\sigma_\theta$: SS friction volatility & $\log(0.008)$ & Guerr\'on-Quintana and Jinnai (2019b) \\
\bottomrule
\end{tabular}
\end{table}

The volatility innovation scale $\nu_j$ is calibrated so that a one-standard-deviation
positive uncertainty innovation $u_{j,t}=1$ doubles the level of $\sigma_{j,t}$
relative to its SS:
\begin{equation}
\nu_j = \frac{\log 2}{(1-\rho_{\sigma_j}^2)^{1/2}}.
\tag{Nu-FVGQ}
\end{equation}
With $\rho_{\sigma_j}=0.75$ this gives $\nu_j \simeq 1.048$, which is the value
used in the codebase.

The deterministic SS values that underpin §\ref{sec:fvgq-closed-form} and
§\ref{sec:third-order-projection} follow from this calibration. In particular,
$\bar R = 1/\beta \simeq 1.0060$, $\bar R^k \simeq 0.0273$, and the BCA-Euler
forward coefficient is
\[
\omega_\tau = \frac{1-\delta}{\bar R^k + 1 - \delta} \simeq 0.973.
\]

\medskip

\noindent\textbf{Bibliographic note.} References in this section --- Com\'in
and Gertler (2006), Cooley and Prescott (1995), Jermann and Quadrini (2012),
Fern\'andez-Villaverde et al.\ (2015a), and Guerr\'on-Quintana and Jinnai
(2019b) --- appear inline without a formal bibliography. A future revision
will compile a \texttt{.bib} file. For the present draft they should be read
as pointers to the standard literature listed in the FVGQ paper.

\section{Generic BCA Wedge Projection}
\label{sec:generic-projection}

\subsection{BCA Prototype Euler Equation}

The representative-household Euler equation in a CKM-style BCA prototype can be
written, after first-order log linearization around the deterministic steady
state, as
\begin{equation}
\hat\tau_{x,t}-\gamma\tilde C_t
=
\E_t\!\left[
-\gamma\tilde C_{t+1}
+\omega_K\tilde R^k_{t+1}
+\omega_\tau\hat\tau_{x,t+1}
\right].
\tag{BCA-E}
\end{equation}
Here $\gamma$ is relative risk aversion, $\tilde C_t$ and
$\tilde R^k_{t+1}$ are the observed consumption and capital-return objects
fed into the BCA prototype, and $\omega_K,\omega_\tau$ are steady-state
coefficients.

Suppose that the true data-generating process contains stochastic volatility.
Its Euler equation, expanded to second order, can be written as
\begin{equation}
-\gamma\tilde C_t
=
\E_t\!\left[
-\gamma\tilde C_{t+1}
+\omega_K\tilde R^k_{t+1}
\right]
+\frac{1}{2}\Phi(\Omega_t)
+o(\|\Delta\Omega_t\|).
\tag{DGP-E}
\end{equation}
The term $\Phi(\Omega_t)$ collects risk corrections generated by conditional
variances, conditional covariances, and Hessian terms.

Feeding the data generated by (DGP-E) into the BCA prototype and comparing
(BCA-E) with (DGP-E) gives the projection equation
\begin{equation}
\boxed{
\hat\tau_{x,t}
=
\omega_\tau\E_t\hat\tau_{x,t+1}
+\frac{1}{2}\Phi(\Omega_t)
}.
\tag{P}
\end{equation}
This equation is the key object of the report.

\subsection{Closed-Form Projection}

\begin{assumption}[Stochastic-volatility state]
The stochastic-volatility state satisfies
\[
\Delta\Omega_{t+1}
=
R_\Omega\Delta\Omega_t+\eta_{t+1},
\qquad
\rho(R_\Omega)<1,
\]
where $\rho(R_\Omega)$ is the spectral radius.
\end{assumption}

\begin{assumption}[Local differentiability]
$\Phi$ is differentiable near $\bar\Omega$ and satisfies
\[
\Phi(\Omega_t)
=
\Phi(\bar\Omega)
+\bm{\phi}'\Delta\Omega_t
+o(\|\Delta\Omega_t\|).
\]
\end{assumption}

\begin{theorem}[Wedge Projection]
Under the assumptions above and $|\omega_\tau|<1$, the locally bounded solution
to the projection equation (P) is
\begin{equation}
\boxed{
\hat\tau_{x,t}
=
\frac{\Phi(\bar\Omega)}{2(1-\omega_\tau)}
+
\frac{1}{2}\bm{\phi}'
(I-\omega_\tau R_\Omega)^{-1}
\Delta\Omega_t
+o(\|\Delta\Omega_t\|)
}.
\tag{T1}
\end{equation}
If $\Omega_t$ is a scalar AR(1) process,
$\Delta\Omega_{t+1}=\rho_\Omega\Delta\Omega_t+\eta_{t+1}$, then
\begin{equation}
\hat\tau_{x,t}
=
\frac{\Phi(\bar\Omega)}{2(1-\omega_\tau)}
+
\underbrace{
\frac{\phi}{2(1-\omega_\tau\rho_\Omega)}
}_{\kappa_x}
(\Omega_t-\bar\Omega)
+o(|\Omega_t-\bar\Omega|).
\tag{T1s}
\end{equation}
\end{theorem}

\begin{proof}
Forward iteration of (P) gives
\[
\hat\tau_{x,t}
=
\frac{1}{2}\sum_{j=0}^{\infty}\omega_\tau^j\E_t\Phi(\Omega_{t+j}).
\]
The local approximation implies
\[
\E_t\Phi(\Omega_{t+j})
=
\Phi(\bar\Omega)+\bm{\phi}'R_\Omega^j\Delta\Omega_t+o(\|\Delta\Omega_t\|).
\]
Substitution and the geometric-series identities
\[
\sum_{j=0}^{\infty}\omega_\tau^j=\frac{1}{1-\omega_\tau},
\qquad
\sum_{j=0}^{\infty}(\omega_\tau R_\Omega)^j
=(I-\omega_\tau R_\Omega)^{-1}
\]
yield (T1). The scalar formula is the one-dimensional special case.
\end{proof}

The theorem has three immediate implications.
\begin{itemize}
  \item \textbf{Persistence pass-through.} BCA may recover a wedge that looks
  like a persistent AR(1) first-order wedge, even though the underlying force is
  a stochastic-volatility state.
  \item \textbf{Amplification.} If $\omega_\tau\rho_\Omega$ is close to one,
  the second-order risk correction is amplified by
  $(1-\omega_\tau\rho_\Omega)^{-1}$.
  \item \textbf{Sign discipline.} The sign of the investment-wedge response is
  governed by $\phi=\Phi'(\bar\Omega)$. This makes it possible to distinguish
  second-order Euler channels from first-order static-FOC channels.
\end{itemize}

\section{FVGQ Closed Form for \texorpdfstring{$\Phi(\Omega_t)$}{Phi(Omega)}}
\label{sec:fvgq-closed-form}

\subsection{Entrepreneur Euler Equation}

This section specializes the generic projection equation (P) of
§\ref{sec:generic-projection} to the FVGQ environment described in
§\ref{sec:fvgq-model}. To keep the algebra transparent, we work in the
representative-entrepreneur, always-binding-collateral, no-adjustment-cost
limit of (F-FVGQ). The two-agent extension is restored in
§\ref{sec:two-agent}.

Let $\lambda_t=u'(c_t)$ be the entrepreneur's marginal utility of
consumption. The bond and capital first-order conditions are
\begin{align}
\lambda_t &= \beta R_t\E_t\lambda_{t+1}+\mu_t, \tag{F1}\\
\lambda_t &= \beta\E_t[\lambda_{t+1}(R^k_{t+1}+1-\delta)]
+\mu_t\theta_t(1-\delta). \tag{F2}
\end{align}
Eliminating the collateral multiplier $\mu_t$ gives
\begin{equation}
\lambda_t[1-\theta_t(1-\delta)]
=
\beta\E_t[\lambda_{t+1}(R^k_{t+1}+1-\delta)]
-\beta R_t\theta_t(1-\delta)\E_t\lambda_{t+1}.
\tag{E-Euler}
\end{equation}
This Euler equation is the source of the FVGQ risk-correction terms.

\subsection{Single-Expectation Reformulation}
\label{subsec:single-expectation}

A na\"ive cumulant expansion that handles the two right-hand side
expectations of (E-Euler) separately and then combines them linearly is
incorrect, because $\log[\,w_R B^R - w_\theta B^\theta\,]$ is not a linear
function of $\log B^R$ and $\log B^\theta$. The correct procedure first
collapses the right-hand side into a single conditional expectation and only
then applies the cumulant identity.

Treating the bond rate as predetermined at the deterministic SS,
$R_t \simeq \bar R$, the right-hand side of (E-Euler) becomes
\begin{equation}
\beta \E_t G_{t+1},
\qquad
G_{t+1} \;\equiv\;
\lambda_{t+1} \big[R^k_{t+1} + 1-\delta - \bar R\,\theta_t(1-\delta)\big].
\tag{G-FVGQ}
\end{equation}
At the deterministic SS, $\bar G = \bar\lambda \bar X$ with
\[
\bar X \;\equiv\; \bar R^k + 1-\delta - \bar R \bar\theta(1-\delta)
\;=\; \frac{1-\bar\theta(1-\delta)}{\beta}.
\]
Defining log deviations $\hat g_{t+1} \equiv \log(G_{t+1}/\bar G)$ and
$\hat r^k_{t+1} \equiv \log[(R^k_{t+1}+1-\delta)/(\bar R^k+1-\delta)]$, a
first-order linearization gives
\begin{equation}
\hat g_{t+1}
\;=\;
\hat\lambda_{t+1} + \omega_K \hat r^k_{t+1} - \omega_\theta \hat\theta_t,
\qquad
\omega_K \;\equiv\; \frac{\bar R^k+1-\delta}{\bar X},
\qquad
\omega_\theta \;\equiv\; \frac{\bar R\bar\theta(1-\delta)}{\bar X}.
\tag{$\hat g$-FVGQ}
\end{equation}
The SS condition $\bar X = \bar R^k+1-\delta - \bar R\bar\theta(1-\delta)$
forces the identity
\[
\omega_K - \omega_\theta = 1.
\]
Numerically (Table~\ref{tab:fvgq-calib}), $\omega_K = 1.1099$ and
$\omega_\theta = 0.1099$.

\subsection{Risk-Correction Decomposition}

Take logs of (E-Euler) using the single-expectation form (G-FVGQ) and apply
the standard cumulant identity to $\log\E_t e^{\hat g_{t+1}}$:
\begin{equation}
\hat\lambda_t
\;=\;
\E_t \hat\lambda_{t+1} + \omega_K \E_t \hat r^k_{t+1}
- \omega_\theta \hat\theta_t
+ \tfrac{1}{2}\Phi(\Omega_t) + o(\sigma^2),
\tag{DGP-E-FVGQ}
\end{equation}
where the second-order risk correction is
\begin{equation}
\boxed{
\Phi(\Omega_t) \;=\; \Var_t(\hat g_{t+1})
\;=\;
\gamma^2 \Sigma_{cc}(\Omega_t)
- 2\gamma\,\omega_K\,\Sigma_{cR}(\Omega_t)
+ \omega_K^2\,\Sigma_{rr}(\Omega_t)},
\tag{FVGQ-$\Phi$}
\end{equation}
with the standard CRRA substitution $\hat\lambda_{t+1}=-\gamma\hat c_{t+1}$
and
\[
\Sigma_{cc}(\Omega_t)\equiv\Var_t(\hat c_{t+1}),
\qquad
\Sigma_{cR}(\Omega_t)\equiv\Cov_t(\hat c_{t+1},\hat r^k_{t+1}),
\qquad
\Sigma_{rr}(\Omega_t)\equiv\Var_t(\hat r^k_{t+1}).
\]
The decomposition into named channels is therefore
\begin{align}
\Phi^{\mathrm{prec}}(\Omega_t)
&= \gamma^2 \,\Sigma_{cc}(\Omega_t),
\tag{Prec}\\
\Phi^{\mathrm{rp}}(\Omega_t)
&= -2\gamma\,\omega_K\,\Sigma_{cR}(\Omega_t),
\tag{RP}\\
\Phi^{\mathrm{rr}}(\Omega_t)
&= \omega_K^2\,\Sigma_{rr}(\Omega_t).
\tag{RR}
\end{align}

\paragraph{Comparison with the previous draft.}
An earlier version of this report contained
$\Phi^{\mathrm{prec}}=\tfrac{1}{2}\gamma^2[1-\bar\theta(1-\delta)]\Sigma_{cc}$
and $\Phi^{\mathrm{rp}}=-\gamma\beta(\bar R^k+1-\delta)\Sigma_{cR}$, derived
by applying the cumulant identity to each expectation in (E-Euler) separately
and then combining linearly. That step is invalid, and the resulting formulas
are uniformly suppressed by a factor of $\tfrac{1}{2}[1-\bar\theta(1-\delta)]
\simeq 0.453$. The corrected formulas above are roughly $2.21$ times larger.
The earlier version also lacked the $\Phi^{\mathrm{rr}}$ term, which is a
genuine third channel: levered exposure to capital returns when
$\omega_K>1$.

\paragraph{The collateral channel is third order.}
An earlier draft also contained
$\Phi^{\mathrm{coll}}(\Omega_t)
= \beta\bar R\bar\theta(1-\delta)\nu_\theta\sigma_{e,t}^2$,
where $\nu_\theta = \partial \E[\log\lambda]/\partial\sigma_e^2$ is the
sensitivity of the entrepreneur's marginal utility to the financial-friction
volatility state. In a standard pruned third-order solution, this sensitivity
is a third-order policy-function object: in second order, $\hat\lambda_t$ has
no $\sigma_{e,t}^2$ component beyond the constant precautionary level shift
absorbed by the stochastic SS. We therefore defer the collateral channel to
§\ref{sec:third-order-projection} (specifically (Psi-Coll)) and exclude it
from the second-order $\Phi$.

\paragraph{Sign content.}
The precautionary component $\Phi^{\mathrm{prec}}>0$ raises $\hat\tau_{x,t}$
when consumption variance rises with the volatility state. The risk-premium
component $\Phi^{\mathrm{rp}}<0$ when consumption and capital returns
co-move (the standard equity risk premium sign). The return-variance
component $\Phi^{\mathrm{rr}}>0$ is unambiguous and is amplified by the
collateral leverage $\omega_K^2$.

\subsection{Mapping to \texorpdfstring{$\kappa_x$}{kappa x}}

Near the deterministic steady state, write the second-order policy-function
loadings of the conditional moments as
\begin{align}
\Sigma_{cc}(\Omega_t)
&= a_{cc,A}\sigma_{A,t}^2+a_{cc,e}\sigma_{e,t}^2,\\
\Sigma_{cR}(\Omega_t)
&= a_{cR,A}\sigma_{A,t}^2+a_{cR,e}\sigma_{e,t}^2,\\
\Sigma_{rr}(\Omega_t)
&= a_{rr,A}\sigma_{A,t}^2+a_{rr,e}\sigma_{e,t}^2.
\end{align}
With the corrected (FVGQ-$\Phi$), the volatility-loading coefficients
$\phi_j \equiv \partial\Phi/\partial\sigma_{j,t}^2|_{\bar\Omega}$ are
\begin{align}
\phi_A
&= \gamma^2 a_{cc,A}
- 2\gamma\,\omega_K\, a_{cR,A}
+ \omega_K^2\, a_{rr,A},
\tag{phiA}\\
\phi_e
&= \gamma^2 a_{cc,e}
- 2\gamma\,\omega_K\, a_{cR,e}
+ \omega_K^2\, a_{rr,e}.
\tag{phie}
\end{align}
The collateral channel ($\nu_\theta\sigma_{e,t}^2$ contribution) is third
order in the perturbation parameter and appears separately in
§\ref{sec:third-order-projection} as part of $\Psi^{\mathrm{coll}}$.

Theorem 1 gives
\begin{equation}
\hat\tau_{x,t}
= \kappa_{x,A}\Delta\sigma_{A,t}^2 + \kappa_{x,e}\Delta\sigma_{e,t}^2,
\tag{Kappa-map}
\end{equation}
with
\begin{equation}
\kappa_{x,A}
= \frac{\phi_A}{2(1-\omega_\tau\rho_{\sigma_A})},
\qquad
\kappa_{x,e}
= \frac{\phi_e}{2(1-\omega_\tau\rho_{\sigma_e})}.
\tag{Kappa}
\end{equation}

This mapping makes the FVGQ second-order sign predictions testable. Macro
uncertainty has the structure $\phi_A = \gamma^2 a_{cc,A} - 2\gamma\omega_K
a_{cR,A} + \omega_K^2 a_{rr,A}$. Whether $\phi_A>0$ depends on the relative
size of the precautionary $\Sigma_{cc}$ load and the risk-premium covariance
$\Sigma_{cR}$ load. Financial uncertainty has analogous structure in
$(\sigma_{e,t}^2)$, with the collateral contribution arriving at third order
through $\Psi^{\mathrm{coll}}$ and not via $\phi_e$ above.

The third loading $a_{rr,j}$ ($j\in\{A,e\}$) is a new requirement of the
corrected derivation. It quantifies how the conditional variance of capital
returns moves with the volatility state. We extract it from the Dynare policy
function in §\ref{sec:numerical-validation} alongside $a_{cc,j}$ and
$a_{cR,j}$.

\section{Two-Agent Risk-Sharing Wedge}
\label{sec:two-agent}

The full FVGQ model has workers and entrepreneurs. The BCA prototype, however,
uses aggregate NIPA consumption $C_t$ as the consumption of a representative
agent. The marginal utility $u'(C_t)$ used by the BCA observer is therefore not
the same as the entrepreneur's marginal utility $u'(c_{E,t})$, which is the
object that prices investment in the true model. This difference creates an
aggregation wedge.

\subsection{Aggregation Identity}

Let
\[
C_t=\omega_W c_{W,t}+\omega_E c_{E,t},
\qquad
\omega_W=0.94,\quad \omega_E=0.06.
\]
Define steady-state consumption shares
\[
s_W=\frac{\omega_W\bar c_W}{\bar C},
\qquad
s_E=\frac{\omega_E\bar c_E}{\bar C},
\qquad
s_W+s_E=1.
\]
A second-order expansion of aggregate consumption gives
\begin{equation}
\hat C_t
=
s_W\hat c_{W,t}+s_E\hat c_{E,t}
+\frac{1}{2}s_Ws_E(\hat c_{E,t}-\hat c_{W,t})^2
+O(3).
\tag{Agg-C}
\end{equation}
Let
\[
\delta c_t\equiv\hat c_{E,t}-\hat c_{W,t}.
\]
Under CRRA preferences,
\[
\log u'(C_t)-\log u'(c_{E,t})
=
-\gamma(\hat C_t-\hat c_{E,t})
=
\gamma s_W\delta c_t
-\frac{1}{2}\gamma s_Ws_E\delta c_t^2
+O(3).
\tag{MU-gap}
\]
The first term is a first-order risk-sharing gap between entrepreneurs and
workers. The second term is a Jensen aggregation term.

\subsection{Projection Contribution}

The aggregation identity (Agg-C) gives
$\hat c_{E,t} = \hat C_t + s_W\delta c_t - \tfrac{1}{2}s_Es_W\delta c_t^2
+ O(3)$. Substitute into the entrepreneur Euler (DGP-E-FVGQ) of
§\ref{subsec:single-expectation} (in the form
$\hat\lambda_{E,t} = \E_t\hat\lambda_{E,t+1} + \omega_K \E_t\hat r^k_{t+1}
+ \tfrac{1}{2}\Phi^E$, suppressing the first-moment $\hat\theta_t$ terms).
Using $\hat\lambda_{E,t}=-\gamma\hat c_{E,t}=
-\gamma\hat C_t -\gamma s_W\delta c_t +\tfrac{1}{2}\gamma s_Es_W\delta c_t^2$,
and matching to the BCA prototype Euler written in $\hat C_t$, the wedge
absorbs both the first-moment risk-sharing gap and the second-moment
aggregation correction:
\begin{equation}
\hat\tau_{x,t}
= \omega_\tau\E_t\hat\tau_{x,t+1}
+ \gamma s_W (\E_t\delta c_{t+1} - \delta c_t)
+ \tfrac{1}{2}\Phi^{\mathrm{tot}}(\Omega_t).
\tag{P-2A}
\end{equation}
The total second-order risk correction decomposes as
\begin{equation}
\Phi^{\mathrm{tot}}
=
\Phi^{\mathrm{prec}}
+\Phi^{\mathrm{rp}}
+\Phi^{\mathrm{rr}}
+\Phi^{\mathrm{agg}},
\tag{Phi-tot}
\end{equation}
where the first three terms are evaluated in BCA observables and the
aggregation term collects the residual.

The conversion of the $\Phi^E$ from entrepreneur observables ($\hat c_{E,t}$,
$\hat r^k_{t+1}$) to BCA observables ($\hat C_t$, $\hat r^k_{t+1}$) uses
$\Var(\hat c_E) = \Var(\hat C) + 2 s_W \Cov(\hat C,\delta c) + s_W^2
\Var(\delta c)$ and
$\Cov(\hat c_E,\hat r^k) = \Cov(\hat C,\hat r^k) + s_W \Cov(\delta c,\hat r^k)$.
Combined with the Jensen aggregation contribution
$\tfrac{1}{2}\gamma s_E s_W [\E_t\delta c_{t+1}^2 - \delta c_t^2]$
from the second-order expansion of $\hat C$, the $\Omega_t$-dependent residual
is
\begin{equation}
\boxed{
\Phi^{\mathrm{agg}}(\Omega_t)
\;=\;
2\gamma^2 s_W \Cov_t(\hat C_{t+1},\delta c_{t+1})
+ \gamma^2 s_W^2 \Var_t(\delta c_{t+1})
- 2\gamma\,\omega_K\, s_W \Cov_t(\delta c_{t+1},\hat r^k_{t+1})
+ \gamma s_E s_W \Var_t(\delta c_{t+1})}.
\tag{Agg-$\Phi$}
\end{equation}

\paragraph{Comparison with the previous draft.}
An earlier version of this report wrote
$\Phi^{\mathrm{agg}}=\gamma s_W s_E[\gamma\Var_t(\delta c_{t+1})
-\Cov_t(\delta c_{t+1},\hat r^k_{t+1})]$.
That formula is dimensionally inconsistent: it scales the
risk-premium-aggregation channel by $\gamma s_E s_W \approx 0.054$,
whereas the correct scaling on $\Cov_t(\delta c,\hat r^k)$ is
$2\gamma\omega_K s_W \approx 2.16$. The previous formula thus understated the
aggregation covariance term by a factor of approximately $40$. The leading
$\Var(\delta c)$ coefficient also changes: from
$\gamma^2 s_W s_E \approx 0.054 \gamma$ to a sum
$\gamma^2 s_W^2 + \gamma s_E s_W \approx 0.972\gamma + 0.027\gamma$ (the
covariance with $\hat C$ also enters separately). The economic content remains
that the BCA prototype, by using composite $\hat C$ rather than entrepreneur
$\hat c_E$, mis-prices the second-moment risk faced by the marginal investor.

\paragraph{Note on the $\Cov(\hat C,\delta c)$ term.}
This first cross-covariance has no analogue in the previous draft because the
draft worked from the "$\delta c^2$" Jensen term in isolation. The proper
single-expectation derivation forces the $\hat C$-with-$\delta c$ cross-load
to appear, and it is in fact the dominant aggregation channel since
$\hat C$ is itself macroeconomically active.

\subsection{Linear Policy-Function Interpretation}

Let the one-step-ahead innovation loading of consumption dispersion be
\[
\delta c_{t+1}
=
\chi_A\varepsilon_{A,t+1}
+\chi_\theta\varepsilon_{\theta,t+1}
+\text{state terms},
\]
and let the corresponding loading of capital returns be
\[
\hat r^k_{t+1}
=
r_A\varepsilon_{A,t+1}
+r_\theta\varepsilon_{\theta,t+1}
+\text{state terms}.
\]
If the two shocks are orthogonal, then
\begin{align}
\Var_t(\delta c_{t+1})
&=
\chi_A^2\sigma_{A,t}^2+\chi_\theta^2\sigma_{e,t}^2,\\
\Cov_t(\delta c_{t+1},\hat r^k_{t+1})
&=
\chi_A r_A\sigma_{A,t}^2+\chi_\theta r_\theta\sigma_{e,t}^2.
\end{align}
Define additionally the loadings of the aggregate consumption and the
entrepreneur-worker covariance with $\hat C$ via
\[
\Cov_t(\hat C_{t+1},\delta c_{t+1})
= a_{C\delta,A}\sigma_{A,t}^2 + a_{C\delta,e}\sigma_{e,t}^2.
\]
With (Agg-$\Phi$), the marginal aggregation coefficients become
\begin{align}
\phi_A^{\mathrm{agg}}
&= 2\gamma^2 s_W\, a_{C\delta,A}
+ \gamma^2 s_W^2\, \chi_A^2
- 2\gamma\,\omega_K\, s_W\, \chi_A r_A
+ \gamma s_E s_W\, \chi_A^2,
\tag{aggA}\\
\phi_e^{\mathrm{agg}}
&= 2\gamma^2 s_W\, a_{C\delta,e}
+ \gamma^2 s_W^2\, \chi_\theta^2
- 2\gamma\,\omega_K\, s_W\, \chi_\theta r_\theta
+ \gamma s_E s_W\, \chi_\theta^2.
\tag{agge}
\end{align}
Compared to the previous draft, the new corrections are: (i) the
$\Cov(\hat C,\delta c)$ cross-load $a_{C\delta,j}$ enters with coefficient
$2\gamma^2 s_W \approx 1.94$ rather than vanishing, (ii) the
$\Cov(\delta c,\hat r^k)$ load enters with coefficient
$2\gamma\omega_K s_W \approx 2.16$ rather than $\gamma s_W s_E \approx 0.054$,
and (iii) the leading $\Var(\delta c)$ coefficient is
$\gamma^2 s_W^2 + \gamma s_E s_W \approx 0.972\gamma + 0.027\gamma$ rather
than $\gamma s_W s_E \approx 0.054\gamma$.

Financial uncertainty typically makes entrepreneur consumption much more
sensitive to collateral shocks than worker consumption, so
$\Var_t(\delta c_{t+1})$ rises with $\sigma_{e,t}^2$ and the dominant
$\gamma^2 s_W^2 \chi_\theta^2$ term reinforces a positive investment-wedge
projection. Macro uncertainty is again ambiguous because both workers and
entrepreneurs are exposed to TFP shocks, and now the $\Cov(\hat C,\delta c)$
cross-channel adds a third potentially-canceling term to the residual sign.

\section{Working-Capital Counter-Hypothesis}

The FVGQ baseline is a second-order Euler-equation channel. A working-capital
friction is different: it is a first-order static labor-demand channel. It
should therefore project into the labor wedge rather than the investment wedge.

\subsection{Exact Labor-Wedge Projection}

Suppose firms must finance a fraction $\theta_{wc}$ of the wage bill before
production. Let $q_t$ denote the external finance spread. The firm's static
condition is
\begin{equation}
A_tF_{L,t}
=
w_t[1+\theta_{wc}q_t].
\tag{WC-FOC}
\end{equation}
The BCA intratemporal labor condition is
\begin{equation}
-\frac{U_L(C_t,L_t)}{U_C(C_t,L_t)}
=
\Lambda_{\ell,t}^{BCA} A_tF_{L,t}.
\tag{BCA-L}
\end{equation}
The true household condition gives $-U_L/U_C=w_t$. Combining these conditions
gives the exact projection identity
\begin{equation}
\boxed{
\Lambda_{\ell,t}^{BCA}
=
\frac{w_t}{A_tF_{L,t}}
=
\frac{1}{1+\theta_{wc}q_t}
}.
\tag{WC-P}
\end{equation}
Log-linearizing around $\bar q$ yields
\begin{equation}
\hat\Lambda_{\ell,t}^{BCA}
=
-\underbrace{\frac{\theta_{wc}}{1+\theta_{wc}\bar q}}_{\lambda_{wc}>0}
(q_t-\bar q)
+O((q_t-\bar q)^2).
\tag{WC-LIN}
\end{equation}
If the spread satisfies the first-order approximation
\begin{equation}
q_t-\bar q
=
\alpha_\xi\hat\xi_t+\alpha_\Omega(\Omega_t-\bar\Omega)
+o(\|\hat\xi_t,\Omega_t-\bar\Omega\|),
\qquad
\alpha_\xi>0,\quad \alpha_\Omega\ge 0,
\tag{RP-WC}
\end{equation}
then
\begin{equation}
\boxed{
\hat\Lambda_{\ell,t}^{BCA}
=
-\lambda_{wc}\alpha_\xi\hat\xi_t
-\lambda_{wc}\alpha_\Omega(\Omega_t-\bar\Omega)
+O(\hat\xi_t^2)
+O((\Omega_t-\bar\Omega)^2)
}.
\tag{T2}
\end{equation}

\subsection{Economic Content}

If a financial tightening $\hat\xi_t>0$ or higher uncertainty raises the
external finance spread, then
\[
\hat\Lambda_{\ell,t}^{BCA}<0.
\]
This sign does not depend on Hessian terms, risk-premium covariance, or the
numerical size of a simulation. Moreover, if $\hat\xi_t$ is a first-moment
shock, the labor-wedge response is $O(\sigma)$, while a pure
stochastic-volatility Hessian projection into the investment wedge is
$O(\sigma^2)$. The two mechanisms are therefore distinguishable by both wedge
location and local order.

\section{Unified Sign Identification}

\begin{theorem}[Wedge Sign Identification]
Consider two competing mechanisms.

\medskip
\noindent
\textbf{$H_0^{\mathrm{coll}}$: FVGQ-style second-order collateral channel.}
If the true DGP transmits financial frictions mainly through second-order risk
corrections in the Euler equation, then
\[
\E_t[\hat\tau_{x,t}\mid \Delta\sigma_{e,t}^2]
=
\kappa_{x,e}\Delta\sigma_{e,t}^2
+o(\Delta\sigma_{e,t}^2),
\]
where $\kappa_{x,e}$ is given by (Kappa). In the baseline FVGQ collateral
mechanism, the collateral and aggregation components push
$\kappa_{x,e}$ toward a positive sign. The after-tax labor wedge does not have
a signed first-order negative term driven by the financing shock.

\medskip
\noindent
\textbf{$H_1^{\mathrm{wc}}$: wage-bill working-capital channel.}
If the true DGP contains a wage-bill working-capital friction and uncertainty
or financial tightening raises the external finance spread, then
\[
\E_t[\hat\Lambda_{\ell,t}\mid \Delta\Omega_t]
=
-\lambda_{wc}(\alpha_\xi\rho_{\xi\Omega}+\alpha_\Omega)\Delta\Omega_t
+o(\Delta\Omega_t)
<0,
\]
provided that $\alpha_\xi\rho_{\xi\Omega}+\alpha_\Omega>0$. Working capital
does not directly create a first-order $\alpha_\xi\hat\xi_t$ term in the
capital Euler equation.
\end{theorem}

The testable signatures are summarized below.

\begin{center}
\begin{tabular}{p{0.22\textwidth}p{0.22\textwidth}p{0.18\textwidth}p{0.25\textwidth}}
\toprule
Mechanism & Direct condition & Main BCA wedge & Sign pattern \\
\midrule
FVGQ collateral / SV risk & Intertemporal Euler, second order & Investment wedge $\tau_x$ &
$\kappa_{x,e}>0$ in baseline; $\kappa_{x,A}$ parameter-dependent \\
\addlinespace
Two-agent aggregation & BCA uses $C$, while investment uses $c_E$ & Investment wedge $\tau_x$ plus dispersion term &
Financial channel likely reinforces $\tau_x>0$; macro channel ambiguous \\
\addlinespace
Working capital & Static labor demand, first order & After-tax labor wedge $\Lambda_\ell$ &
$\hat\Lambda_{\ell}<0$ when financing spreads rise \\
\bottomrule
\end{tabular}
\end{center}

Therefore, if the data or simulations show that financial uncertainty mainly
lowers $\log(1-\tau_\ell)$ while the investment-wedge signal is weak, the
evidence supports the working-capital channel. If the signal is concentrated in
the investment wedge, especially with a positive response of $\tau_x$ to
financial uncertainty, the pattern is closer to the FVGQ collateral and
aggregation channel.

\section{Third-Order Projection Extension}
\label{sec:third-order-projection}

Theorem~1 captures only the second-order risk correction $\Phi(\Omega_t)$. This
section extends the projection equation to third order in the perturbation
parameter, isolating the new objects that enter the BCA wedge under a pruned
third-order solution of the FVGQ environment of §\ref{sec:fvgq-model}. The main
new result is Theorem~\ref{thm:third-order}: under standard regularity
conditions, the locally bounded solution to the third-order projection equation
combines a refined linear coefficient on $\Delta\Omega_t$, a state-dependent
cross term in $(\hat z_t,\Delta\Omega_t)$, and an explicit asymmetry component
that is generated by the exponential structure of the conditional variance.

\subsection{Why the Second-Order Projection Is Insufficient for FVGQ}

Two observations make a third-order extension necessary in the FVGQ setting.

\medskip

\noindent\textit{The FVGQ benchmark uses pruned third order.} FVGQ Section~6
solves the model with a pruned third-order perturbation around the
deterministic SS. The reason is structural: a second-order policy function
contains no terms that involve only the uncertainty shocks $u_{j,t}$. Such
terms first appear at third order, in cross-products of the form
$\sigma_{j,t-1} u_{j,t}$ and in the constant $\nu_j^2$ that scales the vol-of-vol
correction. Without those third-order terms, the impact response of every
endogenous variable to a $u_{j,t}$ innovation is identically zero.

\medskip

\noindent\textit{Theorem~1 cannot recover impact GIRFs.}
The Theorem~1 projection coefficient $\kappa_x$ measures
how $\hat\tau_{x,t}$ moves with $\Delta\Omega_t$, that is, with the
\emph{state} of the conditional variance. It does not measure how
$\hat\tau_{x,t}$ moves with the \emph{innovation} $u_{j,t}$ that perturbs
$\sigma_{j,t}$ at impact. Mechanically, in second order the law of motion of
the BCA wedge satisfies
\[
\hat\tau_{x,t} - \hat\tau_{x,t-1}
=
\kappa_x \rho_{\sigma_j} \Delta\sigma^2_{j,t-1}
+ \kappa_x \big(\sigma^2_{j,t}-\rho_{\sigma_j}\sigma^2_{j,t-1}\big)
+ \text{(\emph{not} } u_{j,t}\text{)}.
\]
Even though the change $\sigma^2_{j,t}-\rho_{\sigma_j}\sigma^2_{j,t-1}$ is
driven by $u_{j,t}$, that mapping is \emph{linear} only at third order in the
perturbation. At second order, $\sigma^2_{j,t}$ is treated as exogenous state
without internal dynamics linked to $u_{j,t}$. The local-projection regression
$\hat\tau_{x,t}\sim u_{j,t}$ is therefore not identified by Theorem~1 alone.

\medskip

\noindent\textit{Three new objects.} A correctly specified third-order
projection introduces the following three new objects beyond the second-order
$\Phi(\Omega_t)$:
\begin{itemize}
\item $\Psi(\Omega_t)$: pure third-order risk correction, capturing the third
cumulant of the log Euler residual conditional on $\Omega_t$.
\item $\Gamma(\hat z_t,\Omega_t)$: a level-volatility cross term, capturing
how the conditional Hessian depends on the level state $\hat z_t$.
\item Asymmetry contribution: a sign-dependent piece generated by the
exponential parametrization $\sigma^2_{j,t}=\exp(2\sigma_{j,t})$ of the
conditional variance.
\end{itemize}

\subsection{Third-Order Expansion of the True Euler Equation}

Throughout this section we maintain the FVGQ Euler equation derived in
§\ref{sec:fvgq-closed-form}:
\begin{equation}
\lambda_t [1-\theta_t(1-\delta)]
= \beta\E_t G_{t+1},
\qquad
G_{t+1}\equiv \lambda_{t+1}\big[R^k_{t+1}+1-\delta -\bar R\,\theta_t(1-\delta)\big].
\tag{E-Euler-G}
\end{equation}
The right-hand side is now a single conditional expectation, which makes the
log expansion clean. (Note that $R_t=\bar R$ at the deterministic SS by the
representative-bond pricing condition; we treat $R_t$ as predetermined when
computing $\E_t$.) Taking $\log$ on both sides and using the
cumulant-generating-function identity
\[
\log\E_t[Y] = \E_t[\log Y]
+ \tfrac{1}{2}\Var_t(\log Y)
+ \tfrac{1}{6}\kcum_3(\log Y\mid \mc{F}_t)
+ \tfrac{1}{24}\kcum_4(\log Y\mid\mc{F}_t)
+ \cdots,
\]
applied to $Y=G_{t+1}/\bar G$, yields
\begin{equation}
0 = \E_t[m_{t+1}+r^k_{t+1}-\hat\theta_{\mathrm{ad},t}]
+\tfrac{1}{2}\Phi(\Omega_t)
+\tfrac{1}{6}\Psi(\Omega_t)
+\Gamma(\hat z_t,\Omega_t)
+ o(\sigma^3),
\tag{DGP-E3}
\end{equation}
where:
\begin{itemize}
  \item $m_{t+1}\equiv \log\lambda_{t+1}-\log\lambda_t$ is the log SDF
  increment;
  \item $r^k_{t+1}\equiv\log[R^k_{t+1}+1-\delta-\bar R\theta_t(1-\delta)]$ is
  the log adjusted gross return;
  \item $\hat\theta_{\mathrm{ad},t}$ collects the contemporaneous
  $\theta_t$-dependent terms (predetermined at $t$);
  \item $\Phi(\Omega_t)=\Var_t(\log G_{t+1})$ is the second-cumulant
  contribution introduced in (FVGQ-Phi);
  \item $\Psi(\Omega_t)=\kcum_3(\log G_{t+1}\mid\mc{F}_t)$ is the third-cumulant
  contribution; and
  \item $\Gamma(\hat z_t,\Omega_t)$ is a residual cross term capturing the
  state-dependence of $\Phi$ and $\Psi$ themselves.
\end{itemize}

\subsubsection*{Structure of $\Psi$ under Gaussian innovations}

For a vector of innovations $\bs{\varepsilon}_{t+1}\sim\mc{N}(0,I)$, with
$\log G_{t+1}$ admitting the second-order expansion
\[
\log G_{t+1}-\overline{\log G}
= \bs{\xi}_1(\Omega_t,\hat z_t)'\bs{\varepsilon}_{t+1}
+ \tfrac{1}{2}\bs{\varepsilon}_{t+1}'\bs{\Xi}_2(\Omega_t,\hat z_t)\bs{\varepsilon}_{t+1}
+ o(\|\bs{\varepsilon}_{t+1}\|^2),
\]
direct computation of the third cumulant of a Gaussian quadratic form gives
\begin{equation}
\Psi(\Omega_t)
= 8\,\tr\!\big[\big(\bs{\Xi}_2\bs{\Sigma}_\varepsilon(\Omega_t)\big)^3\big]
+ 6\,\bs{\xi}_1'\bs{\Sigma}_\varepsilon(\Omega_t)\bs{\Xi}_2\bs{\Sigma}_\varepsilon(\Omega_t)\bs{\xi}_1
+ o(\sigma^3),
\tag{$\Psi$-Gen}
\end{equation}
where $\bs{\Sigma}_\varepsilon(\Omega_t)$ is the conditional covariance matrix
of $\bs{\varepsilon}_{t+1}$, which is a linear function of the volatility state
$\Omega_t$. The two terms in $(\Psi$-Gen$)$ are the cubic
analogues of the precautionary and risk-premium terms in $\Phi$: the first is
``vol-of-vol amplification'' and the second is a ``state-amplified covariance''
contribution.

\subsubsection*{Structure of $\Gamma$}

If $\bs{\xi}_1$ and $\bs{\Xi}_2$ depend on $\hat z_t$, then $\Phi$ and $\Psi$
also depend on $\hat z_t$ through the third-order tensor of $\log G$. Locally,
\begin{equation}
\Gamma(\hat z_t,\Omega_t)
=
\hat z_t'\,\bs{G}_{z\Omega}\,\Delta\Omega_t
+
\hat z_t'\,\bs{G}_{zz}\,\hat z_t\cdot\bar\Phi'
+
\bs{\xi}_1(\bar\Omega,\bar z)'\frac{\partial\bs{\Sigma}_\varepsilon}{\partial\hat z}
\bigg|_{\bar z}\hat z_t\cdot\Delta\Omega_t
+ o(\sigma^3),
\tag{$\Gamma$-Gen}
\end{equation}
where $\bs{G}_{z\Omega}$ and $\bs{G}_{zz}$ are second-derivative tensors of
$\log G$ at the deterministic SS. The first term, $\hat z_t' \bs{G}_{z\Omega}
\Delta\Omega_t$, is the most important new object: it generates a state-dependent
projection in which the response of $\hat\tau_{x,t}$ to a volatility innovation
depends on the level state. We will write this term simply as $\Gamma_{z\Omega}
\hat z_t\Delta\Omega_t$ in the scalar case below.

\subsection{Statement and Proof of the Third-Order Projection}

\begin{assumption}[Joint linear dynamics]
\label{ass:joint-linear}
The level state $\hat z_t$ and the conditional-variance state $\Omega_t$
satisfy a joint VAR(1):
\[
\begin{pmatrix}\hat z_{t+1}\\ \Delta\Omega_{t+1}\end{pmatrix}
=
\begin{pmatrix}R_z & 0\\ 0 & R_\Omega\end{pmatrix}
\begin{pmatrix}\hat z_t\\ \Delta\Omega_t\end{pmatrix}
+
\begin{pmatrix}\bs{\eta}_{z,t+1}\\ \bs{\eta}_{\Omega,t+1}\end{pmatrix},
\]
with $\rho(R_z),\rho(R_\Omega)<1$ and $\E_t[\bs{\eta}_{z,t+1}\otimes
\bs{\eta}_{\Omega,t+1}]=0$ (innovations to the level state and to the volatility
state are uncorrelated to leading order).
\end{assumption}

\begin{assumption}[Local differentiability of risk corrections]
\label{ass:risk-diff}
$\Phi$ and $\Psi$ are once differentiable near $\bar\Omega$:
\[
\Phi(\Omega_t)=\Phi(\bar\Omega)+\bs{\phi}'\Delta\Omega_t+o(\|\Delta\Omega_t\|),
\qquad
\Psi(\Omega_t)=\Psi(\bar\Omega)+\bs{\psi}'\Delta\Omega_t+o(\|\Delta\Omega_t\|).
\]
The cross term has the bilinear local form
\[
\Gamma(\hat z_t,\Omega_t)=\bar\Gamma+\bar\Gamma_z'\hat z_t+\bar\Gamma_\Omega'
\Delta\Omega_t+\hat z_t'\bs{\Gamma}_{z\Omega}\Delta\Omega_t+o(\sigma^3).
\]
\end{assumption}

\begin{theorem}[Third-Order Wedge Projection]
\label{thm:third-order}
Let Assumptions~\ref{ass:joint-linear} and \ref{ass:risk-diff} hold,
and let $|\omega_\tau|<1$. Then the locally bounded solution to the
third-order projection equation
\[
\hat\tau_{x,t}=\omega_\tau\E_t\hat\tau_{x,t+1}
+\tfrac{1}{2}\Phi(\Omega_t)+\tfrac{1}{6}\Psi(\Omega_t)+\Gamma(\hat z_t,\Omega_t)
\]
admits the affine-bilinear form
\begin{equation}
\boxed{
\hat\tau_{x,t}
= \bar\tau_x
+ \bs{\kxtwo}{}'\Delta\Omega_t
+ \bs{\kxvov}{}'\Delta\Omega_t
+ \bar{\bs{\kappa}}_z'\hat z_t
+ \hat z_t'\bs{\kxcross}\Delta\Omega_t
+ o(\sigma^3)},
\tag{T2}
\end{equation}
where
\begin{align}
\bar\tau_x &=
\frac{\Phi(\bar\Omega)/2+\Psi(\bar\Omega)/6+\bar\Gamma}{1-\omega_\tau},
\tag{T2a}\\
\bs{\kxtwo} &= \tfrac{1}{2}(I-\omega_\tau R_\Omega)^{-1}\bs{\phi},
\tag{T2b}\\
\bs{\kxvov} &= \tfrac{1}{6}(I-\omega_\tau R_\Omega)^{-1}\bs{\psi}+(I-\omega_\tau
R_\Omega)^{-1}\bar\Gamma_\Omega,
\tag{T2c}\\
\bar{\bs{\kappa}}_z &= (I-\omega_\tau R_z)^{-1}\bar\Gamma_z,
\tag{T2d}\\
\bs{\kxcross} &= \big(I_z\otimes I_\Omega - \omega_\tau (R_z\otimes
R_\Omega)\big)^{-1}\bs{\Gamma}_{z\Omega}.
\tag{T2e}
\end{align}
\end{theorem}

\begin{proof}
Forward iterating the projection equation as in Theorem~1,
\[
\hat\tau_{x,t}
=\sum_{j=0}^{\infty}\omega_\tau^j\E_t\!\left[\tfrac{1}{2}\Phi(\Omega_{t+j})
+\tfrac{1}{6}\Psi(\Omega_{t+j})+\Gamma(\hat z_{t+j},\Omega_{t+j})\right].
\]
Substitute the local approximations in
Assumption~\ref{ass:risk-diff} and use Assumption~\ref{ass:joint-linear} to
compute conditional means:
\[
\E_t[\Delta\Omega_{t+j}]=R_\Omega^j\Delta\Omega_t,
\qquad
\E_t[\hat z_{t+j}]=R_z^j\hat z_t,
\qquad
\E_t[\hat z_{t+j}\otimes\Delta\Omega_{t+j}]=(R_z^j\otimes R_\Omega^j)(\hat
z_t\otimes\Delta\Omega_t),
\]
where the third equality uses the uncorrelated-innovation condition. Sum the
geometric series term by term:
\begin{align*}
\sum_{j=0}^\infty\omega_\tau^j R_\Omega^j &= (I-\omega_\tau R_\Omega)^{-1},\\
\sum_{j=0}^\infty\omega_\tau^j R_z^j &= (I-\omega_\tau R_z)^{-1},\\
\sum_{j=0}^\infty\omega_\tau^j (R_z^j\otimes R_\Omega^j)
&= (I_z\otimes I_\Omega-\omega_\tau (R_z\otimes R_\Omega))^{-1}.
\end{align*}
The constants $\Phi(\bar\Omega)/2+\Psi(\bar\Omega)/6+\bar\Gamma$ aggregate
into $\bar\tau_x = (\Phi(\bar\Omega)/2+\Psi(\bar\Omega)/6+\bar\Gamma)/(1-\omega_\tau)$.
The remaining identifications yield (T2a)--(T2e).
\end{proof}

\noindent\textbf{Scalar special case.} Suppose $\hat z_t$ and $\Omega_t$ are
scalar AR(1) processes with persistences $\rho_z,\rho_\Omega\in[0,1)$ and that
the cross-derivative is the scalar $\Gamma_{z\Omega}$. Then (T2) simplifies to
\begin{equation}
\hat\tau_{x,t}
= \bar\tau_x
+ \kxtwo \Delta\Omega_t
+ \kxvov \Delta\Omega_t
+ \bar\kappa_z\hat z_t
+ \kxcross \hat z_t\Delta\Omega_t
+ o(\sigma^3),
\tag{T2s}
\end{equation}
with $\kxtwo=\phi/[2(1-\omega_\tau\rho_\Omega)]$,
$\kxvov=\psi/[6(1-\omega_\tau\rho_\Omega)]+\bar\Gamma_\Omega/(1-\omega_\tau
\rho_\Omega)$, $\bar\kappa_z=\bar\Gamma_z/(1-\omega_\tau\rho_z)$, and
\[
\kxcross
= \frac{\Gamma_{z\Omega}}{1-\omega_\tau \rho_z\rho_\Omega}.
\]
Compared to Theorem~1, the persistence factor in $\kxcross$ is
$1-\omega_\tau\rho_z\rho_\Omega$ rather than $1-\omega_\tau\rho_\Omega$. With
$\omega_\tau=0.973$, $\rho_z\approx 0.95$ (the standard FVGQ level
persistence), and $\rho_\Omega=0.75$, this gives
\[
\frac{1}{1-\omega_\tau\rho_\Omega}\simeq 3.70,
\qquad
\frac{1}{1-\omega_\tau\rho_z\rho_\Omega}\simeq 3.26.
\]
The cross-term coefficient is amplified less than the linear coefficient, but
both are nontrivial.

\subsection{FVGQ-Specific Third-Order Decomposition}

We now apply Theorem~\ref{thm:third-order} to the FVGQ Euler equation.
Continuing the decomposition (FVGQ-Phi) into precautionary, risk-premium,
collateral, and aggregation components, the third-cumulant correction has the
analogous decomposition
\begin{equation}
\Psi^{\mathrm{FVGQ}}(\Omega_t)
=
\Psi^{\mathrm{prec}}(\Omega_t)
+\Psi^{\mathrm{rp}}(\Omega_t)
+\Psi^{\mathrm{coll}}(\Omega_t)
+\Psi^{\mathrm{agg}}(\Omega_t).
\tag{FVGQ-Psi}
\end{equation}

\subsubsection*{Precautionary cubic ($\Psi^{\mathrm{prec}}$)}

Using $\hat\lambda_{t+1}=-\gamma\hat c_{t+1}-\tfrac{1}{2}\gamma(\gamma+1)
\Var_t(\hat c_{t+1}/\sigma)+o(\sigma^2)$ from the second-order expansion of
$u'(c)$ under CRRA, the third-cumulant analogue of $\Phi^{\mathrm{prec}}$ is
\begin{equation}
\Psi^{\mathrm{prec}}(\Omega_t)
=
-\gamma^3[1-\bar\theta(1-\delta)]\,\E_t[\hat c_{t+1}^3]
+ 3\gamma^2(\gamma+1)[1-\bar\theta(1-\delta)]\,\E_t[\hat c_{t+1}\Var_t(\hat
c_{t+2}\mid \mc{F}_{t+1})].
\tag{Psi-Prec}
\end{equation}
The first term vanishes for Gaussian innovations of $\hat c_{t+1}$ at first
order in the perturbation parameter, but it is nonzero at the relevant third
order because $\hat c_{t+1}$ contains second-order quadratic terms that
generate non-vanishing third moments. The second term is the
$\E[\hat c \cdot \Var(\hat c)]$-type correction familiar from the
recursive-utility literature.

\subsubsection*{Risk-premium cubic ($\Psi^{\mathrm{rp}}$)}

\begin{equation}
\Psi^{\mathrm{rp}}(\Omega_t)
=
3\gamma\beta(\bar R^k+1-\delta)\,\E_t[\hat c_{t+1}^2\hat r^k_{t+1}]
- 3\gamma^2\beta(\bar R^k+1-\delta)\,\E_t[\hat c_{t+1}(\hat r^k_{t+1})^2]
+ \dots
\tag{Psi-RP}
\end{equation}
Both terms involve mixed moments of consumption and capital returns. They
generalize the second-order $\Sigma_{cR}$ to higher-order coskewness.

\subsubsection*{Collateral channel ($\Psi^{\mathrm{coll}}$)}

The financial-friction sensitivity $\nuth$, defined as the
third-order policy-function loading of the entrepreneur's log marginal utility
on the volatility state,
\[
\nuth \;\equiv\; \frac{\partial \E[\log\lambda_t]}{\partial \sigma_{e,t}^2},
\]
is genuinely a third-order object: in the second-order expansion of
§\ref{sec:fvgq-closed-form} (specifically (FVGQ-$\Phi$)), $\sigma_{e,t}^2$
enters only through $\Sigma_{cc}, \Sigma_{cR}, \Sigma_{rr}$ as conditional
moments of innovations, not through level shifts of $\E[\log\lambda]$. The
collateral channel therefore belongs in this section, not in
§\ref{sec:fvgq-closed-form}.

Decompose the third-order entrepreneur log marginal utility as
\[
\log\lambda_t = (\text{linear in }\hat z_t)
+ \tfrac{1}{2}\nuth\,\sigma_{e,t}^2
+ \tfrac{1}{6}\nuthIII\,\sigma_{e,t}^3
+ \nuthx\,\hat z_t\sigma_{e,t}^2
+ o(\sigma^3),
\]
where $\nuthIII$ is the third-order vol-of-vol sensitivity and
$\nuthx$ is the cross sensitivity to $(\hat z_t,\sigma_{e,t}^2)$. Plugging
into the entrepreneur Euler equation in single-expectation form and
balancing the $\sigma_{e,t}^2$ coefficients on both sides
($\sigma_{e,t}^2$ on the left, $\rho_{\sigma_e}\sigma_{e,t}^2$ on the right
through $\E_t\sigma_{e,t+1}^2 = \rho_{\sigma_e}\sigma_{e,t}^2 +
\text{const}$), the residual that the wedge must absorb is
\begin{equation}
\Psi^{\mathrm{coll}}(\Omega_t)
\;=\;
2\,(1-\rho_{\sigma_e})\,\nuth\,\sigma_{e,t}^2
\;+\;
\tfrac{1}{6}\nuthIII\,\sigma_{e,t}^3,
\tag{Psi-Coll}
\end{equation}
\begin{equation}
\Gamma^{\mathrm{coll}}(\hat z_t,\Omega_t)
\;=\;
\nuthx\,\hat z_t\,\sigma_{e,t}^2.
\tag{Gam-Coll}
\end{equation}
The factor $2(1-\rho_{\sigma_e})$ on the leading $\nuth\sigma_e^2$ term
reflects that $\Psi$ enters the projection multiplied by $\tfrac{1}{6}$ while
the cumulant comes in at order $\tfrac{1}{2}$, so the bookkeeping yields a
factor of $2$ relative to a $\Phi$-style entry; the $(1-\rho_{\sigma_e})$
factor reflects the persistence of the volatility state. Numerically with
$\rho_{\sigma_e}=0.75$, the leading collateral coefficient is
$0.50\,\nuth\sigma_{e,t}^2$.

\paragraph{Comparison with the previous draft.}
An earlier version of this report wrote
$\Phi^{\mathrm{coll}}=\beta\bar R\bar\theta(1-\delta)\nuth\sigma_{e,t}^2$
inside the second-order $\Phi$, with multiplier
$\beta\bar R\bar\theta(1-\delta)\simeq 0.0938$. Both the placement
(second order rather than third) and the multiplier
($\beta\bar R\bar\theta(1-\delta)$ rather than $2(1-\rho_{\sigma_e})\simeq 0.50$)
were incorrect. The corrected coefficient is approximately $5.3$ times the
earlier formula's, while now belonging to $\Psi$ in the third-order
projection.

\subsubsection*{Aggregation cubic ($\Psi^{\mathrm{agg}}$)}

Recall from §\ref{sec:two-agent} that the second-order aggregation
correction is the (Agg-$\Phi$) expression with leading variance coefficient
$\gamma^2 s_W^2$ and a $-2\gamma\omega_K s_W$ covariance term. The
third-order analogue extends each second-cumulant object to its third
cumulant. Schematically,
\begin{equation}
\Psi^{\mathrm{agg}}(\Omega_t)
=
\gamma s_Ws_E
\left[
\gamma^2\,\E_t[\delta c_{t+1}^3]
- 3\gamma\,\E_t[\delta c_{t+1}^2\hat r^k_{t+1}]
+ \E_t[\delta c_{t+1}(\hat r^k_{t+1})^2]
\right].
\tag{Psi-Agg}
\end{equation}
Under independent Gaussian innovations to the policy functions of $\delta c$
and $\hat r^k$, all of these third moments are zero at the second-order policy
function but become non-zero at third order via the same mechanism as in
$\Psi^{\mathrm{prec}}$.

\subsection{Reconciling the Monte Carlo Puzzle of \texorpdfstring{$\nuth<0$}{nu theta < 0}}

The third-order extension provides a structural explanation for the negative
Monte Carlo estimate of $\nuth$ reported in §\ref{subsec:monte-carlo-nu}.
The diagnostic regresses
\[
\E[\log\lambda_t\mid\sigma_{e,0}=\bar\sigma_e + \delta]\;\;\text{on}\;\;
\sigma_e^2 = e^{2(\bar\sigma_e+\delta)},
\]
across an 11-point grid spanning $\bar\sigma_e\pm 2\sigma(\sigma_{e,t})$. With
$\sigma(\sigma_{e,t})=1.05$, the grid extends to $\delta=\pm 2.10$, which
corresponds to $\sigma_e^2$ ranging over almost three orders of magnitude.

Under the third-order expansion above,
\[
\log\lambda \;=\; \mathrm{const} + \tfrac{1}{2}\nuth\sigma_e^2 +
\tfrac{1}{6}\nuthIII\sigma_e^3 + \tfrac{1}{24}\nu_\theta^{(4)}\sigma_e^4 + \dots,
\]
so a regression on $\sigma_e^2$ alone over a wide grid mixes the second-order
coefficient $\nuth$ with the third- and fourth-order coefficients
$\nuthIII,\nu_\theta^{(4)}$. If $\nuthIII<0$ and the grid extends far enough
that $\sigma_e^3$ has a comparable contribution, the OLS slope on $\sigma_e^2$
is biased downward. The numerical observation in §\ref{subsec:monte-carlo-nu}
that the central three grid points give a positive slope (consistent with
small positive $\nuth$) while the wide grid gives a negative slope
($-4.93\times 10^{-2}$) is precisely the signature of higher-order contamination.
A clean estimate of $\nuth$ requires the central-difference window to be
narrow enough that $|\nuthIII\sigma_e/(3\nuth)|\ll 1$.

\subsection{Implications and Three Testable Predictions}

\subsubsection*{Asymmetric impulse responses}

The exponential parametrization $\sigma^2_{e,t}=e^{2\sigma_{e,t}}$ implies
that a positive uncertainty innovation $u_{e,t}=+\delta$ raises $\sigma_e^2$
more than a negative innovation $u_{e,t}=-\delta$ lowers it:
\[
\sigma_e^2(+\delta)-\bar\sigma_e^2 = \bar\sigma_e^2(e^{2\nu_e\delta}-1)
\;>\;
\bar\sigma_e^2(1-e^{-2\nu_e\delta}) = \bar\sigma_e^2-\sigma_e^2(-\delta).
\]
Combined with the linear $\kxtwo$ coefficient, the on-impact wedge response
satisfies
\[
\hat\tau_{x,t}(+\delta)+\hat\tau_{x,t}(-\delta)
\simeq 2\nu_e^2\delta^2\bar\sigma_e^2\kxtwo + O(\delta^3).
\]
A positive $\kxtwo$ therefore predicts that the average of the responses to
$\pm\delta$ uncertainty shocks is strictly positive at impact. This sum can be
recovered from a sign-conditional GIRF or local projection.

\subsubsection*{State-dependent IRFs}

The cross term $\kxcross\hat z_t\Delta\Omega_t$ generates state-dependence in
the impulse response. Specifically, the response of $\hat\tau_{x,t}$ to a unit
$u_{e,t}$ shock is approximately
\[
\frac{\partial\hat\tau_{x,t}}{\partial u_{e,t}}
\simeq 2\nu_e\bar\sigma_e^2\big[\kxtwo+\kxcross\hat z_t\big]
+ O(\sigma^2).
\]
With $\kxcross>0$ (as expected when $\nuthx>0$, which arises naturally if a
collateral-tightening regime amplifies the cross-channel), the wedge response
should be larger when the economy is already in a collateral-tightening
regime ($\hat z_t>0$). This is testable by interacting the LP regressor
$u_{e,t}$ with a measure of collateral conditions at $t$.

\subsubsection*{Vol-of-vol amplification persistence}

In the second-order projection (Theorem~1), the persistence factor on the
wedge response is $1/(1-\omega_\tau\rho_\Omega)\simeq 3.70$. In the third-order
projection, the cross-term factor is $1/(1-\omega_\tau\rho_z\rho_\Omega)\simeq
3.26$, but the pure vol-of-vol contribution within $\kxvov$ inherits the
second-order persistence. A misspecified second-order regression that omits the
cross term and absorbs vol-of-vol into the $\kxtwo$ slope therefore biases the
estimated persistence toward the longer horizon. This is testable by comparing
horizon-by-horizon LP coefficients with and without an explicit
$\hat z_t \cdot u_{e,t}$ interaction.

\section{Numerical Validation from the FVGQ Workspace}
\label{sec:numerical-validation}

The numerical validation uses pruned third-order policy matrices from the FVGQ
preliminary workspace. This section should be read as a diagnostic check of the
closed-form derivation, not as a final calibration exercise.

\subsection{Extracted Coefficients}

The extracted coefficients are listed below. The consumption variable is
\texttt{cspt}, which is composite per-capita consumption. It is not identical
to entrepreneur consumption $c_E$.

\begin{center}
\begin{tabular}{lc}
\toprule
Coefficient & Value \\
\midrule
$a_{cc,A}$ & $4.844\times 10^{-1}$ \\
$a_{cc,e}$ & $4.115\times 10^{-4}$ \\
$a_{cR,A}$ & $1.886\times 10^{-2}$ \\
$a_{cR,e}$ & $-1.786\times 10^{-4}$ \\
$\nu_\theta$ & $0$ at the first/second-order extraction \\
\bottomrule
\end{tabular}
\end{center}

The strongest sanity check is the stochastic-volatility loading:
\[
\frac{\mathrm{ghxu}[c,\mathrm{sigAt},e_A]}{\mathrm{ghu}[c,e_A]}
=0.7500=\rho_{\sigma_A}.
\]
This confirms that the channel by which the volatility state changes the
one-step-ahead conditional variance is present in the Dynare policy function in
exactly the way required by the closed-form derivation.

\subsection{Canonical Calibration Used in the Validation}

The validation uses the actual calibration in
\texttt{fvgq2020\_solveonly.mod}, not the earlier draft values
$\rho_\sigma=0.95$ and $\bar\theta=0.20$.

\begin{center}
\begin{tabular}{lc}
\toprule
Parameter & Value \\
\midrule
$\gamma$ & $2.0$ \\
$\beta$ & $0.994$ \\
$\delta$ & $0.01625$ \\
$\bar\theta$ & $0.0954$ \\
$\bar R^k$ & $0.0273$ \\
$\bar R$ & $1/\beta\simeq 1.0060$ \\
$\rho_{\sigma_A},\rho_{\sigma_e}$ & $0.75$ \\
$\omega_\tau$ & $\simeq 0.973$ \\
$1-\omega_\tau\rho_\sigma$ & $\simeq 0.270$ \\
\bottomrule
\end{tabular}
\end{center}

The amplification factor is therefore about $1/0.270\simeq 3.7$, not the
roughly 15-fold amplification implied by the earlier, incorrect
$\rho_\sigma=0.95$ value.

\subsection{Predicted \texorpdfstring{$\kappa_x$}{kappa} and LP Comparison}

Using the corrected (FVGQ-$\Phi$) formula and the SS-coefficient
$\omega_K=1.1099$ from Table~\ref{tab:fvgq-calib}, the precautionary plus
risk-premium parts of the volatility loadings are
\begin{align*}
\phi_A^{\mathrm{prec+rp}}
&= \gamma^2 a_{cc,A} - 2\gamma\omega_K a_{cR,A}\\
&= 4\cdot 0.4844 - 2\cdot 2 \cdot 1.1099\cdot 0.01886
= +1.854,\\
\phi_e^{\mathrm{prec+rp}}
&= \gamma^2 a_{cc,e} - 2\gamma\omega_K a_{cR,e}\\
&= 4\cdot 4.115\times 10^{-4}
+ 2\cdot 2\cdot 1.1099\cdot 1.786\times 10^{-4}
= +2.439\times 10^{-3}.
\end{align*}
The $\omega_K^2 a_{rr,j}$ return-variance contribution is not yet measured;
it requires extracting the conditional-variance loading of $\hat r^k_{t+1}$
from the Dynare policy-function output and is left as a verification task in
§\ref{sec:research-implications}.

Theorem~1 then yields, for the prec+rp piece alone,
\begin{center}
\begin{tabular}{lcc}
\toprule
Quantity & Previous (incorrect) & Corrected (prec+rp only) \\
\midrule
$\phi_A$ & $+0.840$ & $+1.854$ \\
$\phi_e$ & $+1.10\times 10^{-3}$ & $+2.44\times 10^{-3}$ \\
$\kappa_{x,A}$ & $+1.554$ & $+3.434$ \\
$\kappa_{x,e}$ & $+2.04\times 10^{-3}$ & $+4.51\times 10^{-3}$ \\
\bottomrule
\end{tabular}
\end{center}

Converted into the impact response to a one-unit volatility shock,
\[
u_A\to\hat\tau_x:\quad +2.34\times 10^{-4},
\qquad
u_e\to\hat\tau_x:\quad +4.27\times 10^{-7}.
\]
Both magnitudes are $2.21\times$ larger than under the previous
(incorrect) cumulant-on-difference derivation, and a positive $a_{rr,j}$
loading would push them larger still through the $\omega_K^2\Sigma_{rr}$
channel.

The comparison with preliminary local-projection peaks is
\begin{center}
\begin{tabular}{lccc}
\toprule
Channel & Closed-form sign & LP peak sign & Magnitude note \\
\midrule
$u_A\to\tau_x$ & Positive & Negative at $h=40$ & Sign mismatch at the peak \\
$u_e\to\tau_x$ & Positive & Positive at $h=32$ & Correct sign, magnitude
revisited \\
\bottomrule
\end{tabular}
\end{center}

The macro sign mismatch persists under the correction. The financial-channel
magnitude has roughly doubled and is now more compatible with the LP peak,
but the comparison still depends on (i) whether the LP peak is the right
target (impact-vs-peak issue), and (ii) the magnitude of the unmeasured
$a_{rr,j}$ return-variance loading. This mismatch should not be read as a
failure of the projection theorem; the closed form is an impact and local
projection object, whereas the reported LP number is a long-horizon peak.
The comparison should include $h=0$, the zero crossing, and the peak.

\subsection{Monte Carlo Check for \texorpdfstring{$\nu_\theta$}{nu theta}}
\label{subsec:monte-carlo-nu}

A pruned third-order Monte Carlo check was used to extract the third-order
collateral-risk sensitivity $\nuth$, which under the corrected derivation of
§\ref{sec:fvgq-closed-form} no longer appears in the second-order $\Phi$ but
instead enters $\Psi$ in §\ref{sec:third-order-projection}. The simulation
starts from the stochastic steady state, overwrites the initial
\texttt{siget}, and calls the Dynare decision rule through \texttt{simult\_}.
The grid has 11 points over
\[
\texttt{sige}\pm2\,\mathrm{std}(\texttt{siget}),
\qquad
\texttt{sige}=-4.7975,
\qquad
\mathrm{std}(\texttt{siget})=1.0479,
\]
with 1000 simulations per grid point and $T=200$.

Regressing the first-period cross-sectional Monte Carlo moments on
$\sigma_e^2=\exp(2\,\texttt{siget})$ gives
\begin{center}
\begin{tabular}{lc}
\toprule
Quantity & Global-grid slope \\
\midrule
$\nu_\theta=\partial \E[\log\lambda]/\partial\sigma_e^2$ &
$-4.933\times10^{-2}$ \\
$a_{cc,e}^{MC}$ & $1.635\times10^{-4}$ \\
$a_{cR,e}^{MC}$ & $-1.619\times10^{-4}$ \\
\bottomrule
\end{tabular}
\end{center}

Under the corrected derivation, the collateral channel enters via
$\Psi^{\mathrm{coll}} = 2(1-\rho_{\sigma_e})\nuth\sigma_{e,t}^2 + \tfrac{1}{6}
\nuthIII\sigma_{e,t}^3$ in (Psi-Coll), with leading coefficient
$2(1-\rho_{\sigma_e}) = 0.50$ rather than the previous draft's
$\beta\bar R\bar\theta(1-\delta) = 0.0938$. The implied $\sigma_{e,t}^2$
contribution to the third-order $\hat\tau_x$ projection is
\[
0.50\,\nu_\theta = -2.467\times10^{-2},
\]
and dividing by $6\cdot 2(1-\omega_\tau\rho_{\sigma_e})$ to obtain the
projection coefficient yields a third-order $\kxvov$ contribution of
$-7.61\times 10^{-3}$ (negative, because $\nu_\theta$ is negative on the
wide-grid). For a one-unit \texttt{ue} shock,
$\Delta\sigma_e^2=9.44\times10^{-5}$, so the impact correction is only
$-7.18\times10^{-7}$. Compared to the corrected second-order
$\kappa_{x,e}^{\mathrm{prec+rp}} = 4.51\times 10^{-3}$, this third-order
contribution is small but no longer negligible.

Thus the current Monte Carlo estimate of $\nu_\theta$ is small and negative
on the wide grid, with $\Psi^{\mathrm{coll}}$ contributing in the wrong sign
relative to the prec+rp baseline. As argued in
§\ref{sec:third-order-projection}.5, this wide-grid slope is contaminated by
higher-order terms; the central three grid points imply $\nu_\theta\simeq
+1.5\times 10^{-3}$, which would deliver a small positive third-order
correction.

This estimate should be interpreted as a diagnostic rather than as the final
local derivative. The global 11-point regression gives
$a_{cc,e}^{MC}=1.635\times10^{-4}$, only 39.7 percent of the policy-function
coefficient $4.115\times10^{-4}$. By contrast, the central three grid points
give $4.347\times10^{-4}$, only 5.6 percent above the policy-function value.
The local perturbation check therefore passes near the steady state, while the
wide grid is affected by nonlinearities.

\subsection{Two-Agent Aggregation Correction}

The full aggregation correction was computed from the separate consumption
objects \texttt{ce} and \texttt{cw}. The Dynare policy-function rows were
aligned using \texttt{order\_var}; the first-order aggregation identity
\[
\texttt{cspt}=s_E\,\texttt{ce}+s_W\,\texttt{cw}
\]
then passes as a check. The steady-state consumption shares are
\[
s_W=0.97249,\qquad s_E=0.02751.
\]

The extracted one-step innovation loadings are
\begin{center}
\begin{tabular}{lrrrr}
\toprule
Channel & $\psi_E$ & $\psi_W$ & $\chi$ & $r$ \\
\midrule
$A$ & $-9.852\times10^{-4}$ & $5.812\times10^{-3}$ &
$-8.410\times10^{-1}$ & $2.710\times10^{-2}$ \\
$e$ & $-5.266\times10^{-3}$ & $3.477\times10^{-4}$ &
$-5.892\times10^{-1}$ & $-8.806\times10^{-3}$ \\
\bottomrule
\end{tabular}
\end{center}
These imply
\begin{center}
\begin{tabular}{lccc}
\toprule
Channel & $\partial\Var(\delta c)/\partial\sigma^2$ &
$\partial\Cov(\delta c,r^k)/\partial\sigma^2$ &
$\phi^{\mathrm{agg}}$ \\
\midrule
$A$ & $7.073\times10^{-1}$ & $-2.279\times10^{-2}$ &
$7.692\times10^{-2}$ \\
$e$ & $3.472\times10^{-1}$ & $5.189\times10^{-3}$ &
$3.688\times10^{-2}$ \\
\bottomrule
\end{tabular}
\end{center}

Using the same normalization as the rest of the report,
$\kappa=\phi/[2(1-\omega_\tau\rho_\sigma)]$ and
$2(1-\omega_\tau\rho_\sigma)=0.5405$. Adding the aggregation correction gives
\begin{center}
\begin{tabular}{lrrrr}
\toprule
Channel & $\phi$ original & $\phi$ updated & $\kappa$ updated & LP peak \\
\midrule
$A$ & $8.399\times10^{-1}$ & $9.168\times10^{-1}$ &
$1.696$ & $-1.642\times10^{-3}$ \\
$e$ & $1.105\times10^{-3}$ & $3.799\times10^{-2}$ &
$7.028\times10^{-2}$ & $1.886\times10^{-3}$ \\
\bottomrule
\end{tabular}
\end{center}

The aggregation correction is positive for both channels. It therefore cannot
explain the negative LP peak for $u_A\to\tau_x$; after aggregation,
$\kappa_{x,A}$ remains positive and moves further away from the LP peak sign.
For $u_e$, aggregation substantially raises the predicted coefficient, but it
does not by itself close the full magnitude gap once the comparison is made in
the impact units of the volatility shock.

\section{Research Implications and Next Steps}

The derivation turns the Year 2 BCA methodology from an informal set of
competing hypotheses into a theorem-based framework with testable wedge
signatures. The next steps are:
\begin{enumerate}
  \item \textbf{Fix the numerical comparison.} Extract the LP impulse response
  at $h=0$, at the zero crossing, and at the peak, rather than comparing only
  the peak. The current numerical comparison mixes impact closed-form objects
  with long-horizon LP peaks.

  \item \textbf{Refine the $\nu_\theta$ derivative.} The broad-grid Monte Carlo
  estimate is negative and small, but the central three grid points pass the
  local perturbation sanity check. The next check should use a narrower grid or
  a common-random-number central difference around the stochastic steady state.

  \item \textbf{Reconcile coefficient units.} The aggregation correction
  substantially raises $\kappa_{x,e}$, but direct comparison to LP peaks
  depends on whether the reported object is a coefficient per unit
  $\sigma^2$, an impact response to a unit volatility innovation, or a
  long-horizon projection peak.

  \item \textbf{Implement the working-capital comparison model.} Add wage-bill
  financing to the FVGQ block and verify whether the sign theorem
  $\hat\Lambda_{\ell,t}<0$ holds in simulation.

  \item \textbf{Write the formal appendix.} Complete the second-order expansion
  of (E-Euler), the aggregation wedge derivation, and the working-capital labor
  projection, with a consistent $\Phi$ normalization throughout.

  \item \textbf{Test the three third-order predictions.} The third-order
  projection (Theorem~\ref{thm:third-order}) yields three new identification
  handles. (i) Asymmetric IRFs: estimate sign-conditional GIRFs from the FVGQ
  pruned solution and check whether the response to $u_{e,t}=+\delta$ exceeds
  the response to $u_{e,t}=-\delta$ in absolute value, and whether the average
  of the two responses has the predicted sign of $\kxtwo$. (ii)
  State-dependent IRFs: interact the LP regressor $u_{e,t}$ with a measure of
  $\hat z_t$ (e.g.\ the contemporaneous level of the $\theta_t$ shock or the
  HP-detrended $q_t$) and verify whether the cross coefficient is positive.
  (iii) Vol-of-vol amplification: extract the cubic policy-function coefficient
  $\nuthIII$ from the third-order Dynare output and confirm whether it
  reconciles the wide-grid Monte Carlo regression slope.

  \item \textbf{Sharpen the $\nuthIII$ extraction.} Implement a narrow-grid
  central-difference estimator with common random numbers around the
  stochastic steady state, computing $\nuth$, $\nuthIII$, and $\nuthx$ jointly.
  This isolates the structural source of the wide-grid bias documented in
  §\ref{subsec:monte-carlo-nu}.
\end{enumerate}

Remaining issues that do not undermine tractability include pruning corrections
to the risk-adjusted steady state, identification with correlated
multi-dimensional volatility states, the non-smooth Hessian created by
occasionally binding collateral constraints, and the measurement gap between
composite consumption and entrepreneur marginal utility.

\appendix

\section{Algebraic Details}

\subsection{Forward Solution of the Projection Equation}

Starting from
\[
\hat\tau_{x,t}
=
\omega_\tau\E_t\hat\tau_{x,t+1}
+\frac{1}{2}\Phi(\Omega_t),
\]
repeated substitution gives, for any $J$,
\[
\hat\tau_{x,t}
=
\frac{1}{2}\sum_{j=0}^{J}\omega_\tau^j\E_t\Phi(\Omega_{t+j})
+\omega_\tau^{J+1}\E_t\hat\tau_{x,t+J+1}.
\]
If the solution is locally bounded and $|\omega_\tau|<1$, the last term
converges to zero. Substituting
$\E_t\Delta\Omega_{t+j}=R_\Omega^j\Delta\Omega_t$ yields Theorem 1.

\subsection{Working-Capital Labor Wedge}

The working-capital firm condition and household condition are
\[
A_tF_{L,t}=w_t(1+\theta_{wc}q_t),
\qquad
-\frac{U_L}{U_C}=w_t.
\]
The BCA labor condition requires a wedge such that
\[
w_t=\Lambda_{\ell,t}^{BCA}A_tF_{L,t}.
\]
Therefore
\[
\Lambda_{\ell,t}^{BCA}
=
\frac{1}{1+\theta_{wc}q_t}.
\]
Taking logs and expanding around $\bar q$ gives
\[
\log\Lambda_{\ell,t}^{BCA}
=
-\log(1+\theta_{wc}q_t)
\simeq
-\frac{\theta_{wc}}{1+\theta_{wc}\bar q}(q_t-\bar q).
\]
This is equation (WC-LIN).

\subsection{Two-Agent Consumption Aggregation}

Let $c_{i,t}=\bar c_i\exp(\hat c_{i,t})$. Then
\[
\frac{C_t}{\bar C}
=
s_W\exp(\hat c_{W,t})+s_E\exp(\hat c_{E,t}).
\]
Taking logs and expanding to second order gives
\[
\hat C_t
=
s_W\hat c_{W,t}+s_E\hat c_{E,t}
+\frac{1}{2}
\left[
s_W\hat c_{W,t}^2+s_E\hat c_{E,t}^2
-(s_W\hat c_{W,t}+s_E\hat c_{E,t})^2
\right]
+O(3).
\]
The bracketed term simplifies to
$s_Ws_E(\hat c_{E,t}-\hat c_{W,t})^2$, which proves (Agg-C).

\end{document}
